% !TEX root = ./zprava.tex


% arara: luahbtex: { synctex: on, interaction: nonstopmode, shell: on,options:[--fmt=optex] }

\input ctustyle3
\worktype [O/CZ]
\workname {SFI-semestrální práce}
\faculty {F2}
\department {12135 - ústav výrobních strojů a zařízení}
\title {SFI-semestrální práce}
\author {Ondřej Dohnal}
\supervisor{Ing. Pavel Steinbauer Ph.D.}
\date {Leden 2026}


\makefront
\chap Požadované artefakty semestrálního projektu
\sec Úvodní studie (Vision and scope)

\secc Deklarace záměru: účel, přínos, motivace, provozní prostředí
Projektem je databáze pro sběr dat z laser trackeru AT960, tato databáze slouží k sběru poloh robota pro následné učení evolučních algoritmů.
Výsledek by měl být aplikace která může sbírat data po dlouhé časové intervaly (12 hodin a více) a ukládat je do databáze.

\secc Stakeholdeři(kdo systém používá, kdo ho spravuje, kdo je regulátor/bezpoečnostní autorita)
Uživatel zahajuje měření a na jeho konci dostává databázi s daty, které jdou následně zpracovávat.
Administrátor se stará o chod systému a jeho debugování.

\secc Top-5 klíčových scénářů použití
Měření dat z pohybu robota, ukládání dat v databázi, filtrace dat, export dat

\secc Odhad složitosti a rizik (3-5 hlavních rizik)
Funkčnost přenosu dat, funkčnost měřícího nástroje, ukládání správného datového typu, čitelnost dat

\secc Plán ověřování v jedné větě: "Úspěch je, když.."(měřitelné)
Úspěch je když se dokáže naplnit databáze daty s údaji o poloze koncového efektoru00

\secc Hranice systému(Co je uvnitř a co je už okolí)-Kontext systému(C4 Kontext/kontextový diagram)+seznam událostí/podnětů z okolí
Uvnitř: kód pro spuštění měření, přenos dat do databázového systému, velmi jednoduché UI 
Venku: programování hardware robota a měřícího zařízení, evoluční algoritmy

\secc Seznam rolí a aktérů (včetně bezpečnostních rolí)
Uživatel-spuštění měření
Admin-úprava a správa kódu

\secc Základní omezení: HW rozhraní, standarty, bezpečnostní limity, reálný čas.
HW rozhraní=Leika Laser Tracker AT960 + Aiwa Kukka, reakční čas musí být do 0.5s.

\secc Plán práce týmu + rozdělení rolí a autorizace
Vše dělám sám.


\sec Specifikace požadavků (Requirements Specification)
\secc Use-case model/user stories: hlavní scénáře + alternativní a chybové toky
Jako uživatel chci naměřit kvanta dat pro následné další zpracování.

\secc Seznam funkčních a nefunkčních požadavků s identifikátory (např. FR-01, NFR-03). U každého: popis,priorita,zdroj, verifikace
Funkční systém naplní databázi pozicemi odrazky.
Nefunkční systém bude mít dlouhou odezvu (více jak 5 sekund), data o poloze jsou prázdná

\secc Požadavky na rozhraní: signály, zprávy, API, datové formáty, časování.
Hlavní komunikace mezi rozhraním bude zařizovat ROS2, data se budou přenášet ve formátu float32, frekvence snímání bude do 1 sekundy od změny polohy robota.


\secc Stop-stavy a chování při poruchách - vazba na bezpečnost
Přílišné množství stejných dat --> systém se zastaví a vyhodí chybu
Výpadek DB --> měření se zastaví dokud DB znovu nepoběží

\sec Modely systému (konceptuální návrh)
\secc Doménový model
\secc Dynamický model
\secc Procesní/behaviroální model
\secc Model rozhraní a nasazení

\sec Ověřovací plán
\secc V&V matice
\secc Test strategie po úrovních
\secc Testovací scénáře
\secc Testovací data a očekávané výsledky
\secc Plán záznamu výsledků

\sec Prototyp
\bye